% ---------------------------------------------------------------
%  PHAN I -- MO DAU
% ---------------------------------------------------------------
\section{MỞ ĐẦU}

\subsection{Lí do chọn đề tài}

Kiểm tra, đánh giá là khâu then chốt quyết định chất lượng dạy học. Chương
trình Giáo dục phổ thông 2018 ban hành kèm theo Thông tư số 32/2018/TT-BGDĐT
chuyển trọng tâm đánh giá từ ghi nhớ kiến thức sang đánh giá năng lực; Thông tư
số 22/2021/TT-BGDĐT quy định việc đánh giá học sinh trung học cơ sở và trung
học phổ thông phải bảo đảm tính chính xác, công bằng, vì sự tiến bộ của người
học. Đặc biệt, Quyết định số 764/QĐ-BGDĐT ngày 08/3/2024 của Bộ Giáo dục và Đào
tạo công bố cấu trúc định dạng đề thi tốt nghiệp trung học phổ thông từ năm
2025 với ba dạng câu hỏi trắc nghiệm mới --- nhiều lựa chọn, đúng/sai và trả
lời ngắn --- đã đặt ra yêu cầu mới cho mỗi giáo viên Toán: đề kiểm tra thường
xuyên, định kì ở trường phải sớm làm quen học sinh với cấu trúc đó.

Thực tế tại trường tôi, việc đáp ứng yêu cầu này gặp ba khó khăn lớn. Thứ nhất,
soạn một đề kiểm tra đúng ma trận bốn mức độ với đủ bốn dạng câu hỏi mất từ bốn
đến sáu giờ làm việc, trong khi một giáo viên thường phải soạn nhiều đề cho
nhiều lớp trong cùng một đợt kiểm tra. Thứ hai, các mã đề được tạo bằng cách
đảo thứ tự câu hỏi và phương án nên thực chất vẫn là một đề, học sinh ngồi gần
nhau vẫn có thể trao đổi đáp án. Thứ ba, kho câu hỏi tích luỹ qua nhiều năm nằm
rải rác trong hàng trăm tệp Word và \LaTeX\ không được đặt tên thống nhất, đến
năm sau gần như phải soạn lại từ đầu.

Trong khi đó, Nghị quyết số 57-NQ/TW ngày 22/12/2024 của Bộ Chính trị về đột
phá phát triển khoa học, công nghệ, đổi mới sáng tạo và chuyển đổi số quốc gia
xác định chuyển đổi số là phương thức sản xuất mới, yêu cầu ứng dụng mạnh mẽ
trí tuệ nhân tạo vào mọi lĩnh vực, trong đó có giáo dục. Các mô hình ngôn ngữ
lớn hiện nay đã đủ khả năng hiểu yêu cầu bằng tiếng Việt của giáo viên; tuy
nhiên nếu giao hẳn việc ra đề cho trí tuệ nhân tạo thì lời giải thường sai và
không kiểm chứng được, không thể dùng làm đề kiểm tra chính thức.

Từ ba khó khăn thực tế và một cơ hội công nghệ nêu trên, tôi chọn nghiên cứu và
thực hiện sáng kiến \textit{``\TENSANGKIEN''}, với hướng đi khác các phần mềm
hiện có: trí tuệ nhân tạo chỉ làm nhiệm vụ hiểu yêu cầu của giáo viên, còn việc
sinh câu hỏi, chọn câu theo ma trận và tính điểm đều do mã nguồn đảm nhiệm để
bảo đảm tính chính xác tuyệt đối.

\ghichu{Nếu trường chị có đặc thù riêng (ví dụ lớp đông, thiếu phòng máy, học
sinh vùng khó) thì thêm một hai câu vào đoạn ``Thực tế tại trường tôi'' --- hội
đồng chấm rất thích chi tiết cụ thể của đơn vị.}

\subsection{Mục tiêu của sáng kiến}

\textbf{Mục tiêu chung:} Xây dựng và đưa vào sử dụng một hệ thống ngân hàng đề
môn Toán trên nền web, giúp giáo viên rút ngắn thời gian biên soạn đề kiểm tra,
bảo đảm đề đúng ma trận và mỗi học sinh làm một đề riêng, đồng thời hỗ trợ tổ
chức kiểm tra và chấm điểm trực tuyến.

\textbf{Mục tiêu cụ thể:}

\begin{enumerate}[label=(\arabic*), leftmargin=1.2cm, itemsep=2pt]
  \item Xây dựng chuẩn mã hoá định danh câu hỏi thống nhất theo lớp, chương,
        bài, mức độ nhận thức và dạng câu hỏi, làm cơ sở để máy tự chọn câu.
  \item Xây dựng thư viện hàm sinh câu hỏi bằng Python cho phép tạo ra vô số
        biến thể của cùng một dạng toán, kèm lời giải chi tiết chính xác.
  \item Lập trình thuật toán phân bổ câu hỏi theo ma trận đề, bám số tiết của
        từng bài trong phân phối chương trình.
  \item Tự động sinh đề và lời giải dưới dạng tệp PDF trình bày bằng \LaTeX,
        đạt chuẩn văn bản đề thi.
  \item Triển khai chức năng làm bài, chấm điểm tự động theo thang điểm 10 và
        trả đề, đáp án, điểm số về lớp học Google Classroom của học sinh.
  \item Đánh giá hiệu quả của hệ thống qua thời gian soạn đề của giáo viên và
        kết quả học tập của học sinh.
\end{enumerate}

\subsection{Đối tượng, phạm vi và thời gian thực hiện}

\textbf{Đối tượng nghiên cứu:} Quy trình biên soạn đề kiểm tra và tổ chức kiểm
tra, đánh giá môn Toán ở trường trung học phổ thông; giải pháp tự động hoá quy
trình đó bằng ngân hàng câu hỏi tham số và trí tuệ nhân tạo.

\textbf{Đối tượng áp dụng:} Giáo viên tổ Toán và học sinh các lớp \dotfill\ của
\DONVI.
\ghichu{Điền tên lớp thực nghiệm và lớp đối chứng, ví dụ ``lớp 10A1, 10A2''.}

\textbf{Phạm vi nội dung:} Sáng kiến được triển khai trước hết với chương trình
Toán lớp 10 (bộ sách Kết nối tri thức với cuộc sống), chương \dotfill; phần
khung hệ thống áp dụng được cho mọi khối lớp.
\ghichu{Ghi rõ chương đã hoàn thành ngân hàng câu hỏi, ví dụ ``Chương I. Mệnh
đề và tập hợp''. Không nên ghi rộng hơn phần thực sự đã làm --- hội đồng sẽ
kiểm tra.}

\textbf{Thời gian thực hiện:} Từ tháng \dots\ năm 2025 đến tháng \dots\ năm
2026, gồm ba giai đoạn: nghiên cứu và thiết kế hệ thống; lập trình và xây dựng
ngân hàng câu hỏi; thực nghiệm sư phạm và hoàn thiện.

\subsection{Phương pháp nghiên cứu}

\begin{itemize}[leftmargin=1.2cm, itemsep=2pt]
  \item \textbf{Phương pháp nghiên cứu lí luận:} nghiên cứu Chương trình Giáo
        dục phổ thông 2018, các văn bản chỉ đạo về kiểm tra đánh giá và cấu
        trúc định dạng đề thi, tài liệu về đo lường và đánh giá trong giáo dục.
  \item \textbf{Phương pháp điều tra, khảo sát:} khảo sát giáo viên tổ Toán về
        thời gian và khó khăn khi biên soạn đề; khảo sát học sinh về mức độ
        hứng thú và tính công bằng của hình thức kiểm tra.
  \item \textbf{Phương pháp thực nghiệm sư phạm:} tổ chức dạy học và kiểm tra
        có đối chứng giữa lớp sử dụng hệ thống và lớp không sử dụng.
  \item \textbf{Phương pháp thống kê toán học:} xử lí số liệu điểm kiểm tra để
        so sánh kết quả hai nhóm lớp.
  \item \textbf{Phương pháp chuyên gia:} xin ý kiến tổ chuyên môn và ban giám
        hiệu về chất lượng đề do hệ thống sinh ra.
\end{itemize}

\subsection{Điểm mới của sáng kiến}

So với các giải pháp đã có, sáng kiến có bốn điểm mới:

\begin{enumerate}[label=(\arabic*), leftmargin=1.2cm, itemsep=2pt]
  \item \textbf{Câu hỏi tham số thay cho câu hỏi tĩnh.} Mỗi câu hỏi trong ngân
        hàng là một chương trình sinh câu hỏi, không phải một đoạn văn bản cố
        định; nhờ đó mỗi học sinh nhận một đề khác nhau về số liệu nhưng tương
        đương về dạng toán và độ khó.
  \item \textbf{Phân vai rõ ràng giữa trí tuệ nhân tạo và mã nguồn.} Trí tuệ
        nhân tạo chỉ đọc hiểu yêu cầu bằng tiếng Việt của giáo viên; việc chọn
        câu, tính điểm và tạo đề hoàn toàn do mã nguồn thực hiện nên kết quả
        chính xác và lặp lại được.
  \item \textbf{Ma trận đề đặt trong bảng cấu hình, không gài cứng trong mã.}
        Khi Bộ hoặc Sở thay đổi cấu trúc đề, giáo viên chỉ sửa số liệu trong
        bảng mà không cần lập trình lại.
  \item \textbf{Khép kín một vòng dạy học.} Từ một câu yêu cầu bằng tiếng Việt
        đến đề PDF, lời giải, bài làm trực tuyến, điểm số và sổ điểm Google
        Classroom đều diễn ra trong cùng một hệ thống.
\end{enumerate}

\chenanh[0.9]{01-giao-dien-tong-quan.png}{Giao diện trang chủ hệ thống Ngân hàng đề Toán tại địa chỉ nganhangdechv.tech}

\ghichu{Ảnh 01: chụp toàn màn hình trang chủ sau khi đăng nhập, thấy rõ ô chat
và tên hệ thống. Lưu thành \texttt{skkn/hinhanh/01-giao-dien-tong-quan.png}.
Khi đã có file, ô đỏ sẽ tự biến mất và ảnh hiện ra --- không phải sửa gì trong
bài.}

\clearpage
